\documentclass{article}

\usepackage[utf8]{inputenc}
\usepackage[spanish]{babel}
\usepackage{amssymb}
\usepackage{graphicx}
\usepackage{amsbsy}
\usepackage{pifont}
\usepackage{epsfig}

\begin{document}

\hyphenation{}

\title{Titulo}

\author{Nombre del Alumno}

\maketitle

\section{Introducción}\label{sec:introduccion}
Describir el problema a resolver\cite{ejemplo:2018}
\cite{cialdini1993influence}


\section{Estado del arte}\label{sec:estadoArte}
Trabajos que se han hecho para resolver el problema

\section{Descripción de la propuesta}\label{sec:descripcion}
La propuesta consiste en ....

La idea principal es aplicar ...

\section{Objetivos}\label{sec:objetivos}
\begin{itemize}
 \item \textbf{Objetivo general}
   \begin{itemize}
     \item Desarrollar .....
   \end{itemize}
 \item \textbf{Objetivos Específicos}
   \begin{itemize}
    \item Desarrollar ....
    \item Proponer ...
    \item Implementar ....
    \item Demostrar ....
    \item Evaluar ....
   \end{itemize}
\end{itemize}

\section{Metas}\label{sec:metas}
Las metas de la propuesta se pueden enumerar de la siguiente manera:
\begin{enumerate}
  \item 
\end{enumerate}

\section{Aportación científica esperada}\label{sec:aporte}
Al finalizar este proyecto de investigación se contará con ...
\begin{enumerate}
\item Modelo ...
\item Método ...
\item
\item
\end{enumerate}

\section{Metodología científica}\label{metodologia}
En la etapa inicial, ....\\

Seguidamente, ...


Finalmente, ...


\section{Calendario de actividades}\label{sec:calendario}
El calendario de actividades se divide en las siguientes fases/actividades principales y están detalladas en la Figura xx :\\

Fase 1: Se realizará un estudio de la bibliografía y del estado del arte.\\

Fase 2: .....\\

Fase 3: .....

Aquí introducir diagrama de gantt con el calendario por fases o por actividades.
%\begin{figure}
%\includegraphics[width=13cm, height=7cm]{diagramaGant}
%\caption{Calendario de Actividades}\label{fig:calendario}
%\end{figure}


\bibliographystyle{plain}
\bibliography{protocolo}

\clearpage
% Visto bueno
\vfill
\Large{Vo. Bo.}
\\    \\    \\    \\    \\  \\    \\    \\    \\    \\   
 
\begin{tabular}{@{}p{6cm}p{7cm}@{}}
    \Large{Firma del Alumno:} & \hrulefill \\
                         & \large{Alumno} \par \\
    \\    \\    \\    \\    \\    
    \Large{Directora de tesis:} & \hrulefill \\
                         & \large{Dra. Helena Gómez Adorno} \\
    \Large{} & \\
                         & \large{Investigadora Titular A} \\
    \Large{} &  \\
                         & \large{IIMAS-UNAM} \\
    
\end{tabular}


\end{document}
